% ============================================================
%  Chapter 2 — Background and Literature Review
% ============================================================
\chapter{Background and Literature Review}
\label{cha:background}

The development of retrieval-augmented generation cannot be understood in isolation
from the broader trajectory of language models.
This chapter traces that trajectory, from the limitations of purely parametric models
through the retrieval techniques that underpin modern RAG systems, to the diverse
architectural families that have emerged between 2020 and 2026.
Figure~\ref{fig:rag_timeline} summarises the key milestones.
The chapter closes with an account of RAG evaluation methodology and a review of
related work in clinical decision support.

% ------------------------------------------------------------------
% Figure 2.1 — RAG Evolution Timeline
% ------------------------------------------------------------------
\begin{figure}[htbp]
\centering
\begin{tikzpicture}[
  font=\small,
  year/.style={
    rectangle, rounded corners=3pt,
    fill=gray!20, draw=gray!60,
    minimum width=1.0cm, minimum height=0.55cm,
    align=center, font=\bfseries\small
  },
  above item/.style={
    rectangle, rounded corners=2pt,
    fill=blue!12, draw=blue!35,
    text width=2.7cm, align=center,
    minimum height=0.5cm, font=\footnotesize
  },
  below item/.style={
    rectangle, rounded corners=2pt,
    fill=teal!12, draw=teal!35,
    text width=2.7cm, align=center,
    minimum height=0.5cm, font=\footnotesize
  },
  tick/.style={gray!60, line width=0.4pt}
]

% Main axis
\draw[-{Stealth[length=5pt]}, line width=1.2pt, gray!70]
  (0,0) -- (13.8,0) node[right, font=\small\bfseries, gray!70] {};

% Year positions: 2020=0, 2021=2.3, 2022=4.6, 2023=6.9, 2024=9.2, 2025=11.5, 2026=13.5
\foreach \x/\yr in {0/2020, 2.3/2021, 4.6/2022, 6.9/2023, 9.2/2024, 11.5/2025, 13.5/2026}{
  \draw[tick] (\x, -0.25) -- (\x, 0.25);
  \node[year] at (\x, 0) {\yr};
}

% ABOVE axis items (odd ones, to avoid overlap)
\node[above item] (rag)    at (0,   1.5)  {RAG\\{\tiny Lewis et al.}};
\node[above item] (dpr)    at (0,   2.55) {DPR\\{\tiny Karpukhin et al.}};
\node[above item] (hyde)   at (4.6, 1.5)  {HyDE\\{\tiny Gao et al.}};
\node[above item] (selfrag) at (6.9, 1.5) {Self-RAG\\{\tiny Asai et al.}};
\node[above item] (graphrag) at (9.2,1.5) {GraphRAG\\{\tiny Edge et al.}};
\node[above item] (raptor) at (9.2, 2.55) {RAPTOR\\{\tiny Sarthi et al.}};
\node[above item] (chain)  at (11.5,1.5)  {Chain-of-Retr.\\{\tiny Wang et al.}};
\node[above item] (arag)   at (13.5,1.5)  {A-RAG\\{\tiny Du et al.}};

% BELOW axis items
\node[below item] (colbert)  at (0,  -1.5)  {ColBERT\\{\tiny Khattab \& Zaharia}};
\node[below item] (beir)     at (2.3,-1.5)  {BEIR\\{\tiny Thakur et al.}};
\node[below item] (colbertv2) at (4.6,-1.5) {ColBERTv2\\{\tiny Santhanam et al.}};
\node[below item] (ragas)    at (6.9,-1.5)  {RAGAS\\{\tiny Es et al.}};
\node[below item] (crag)     at (9.2,-1.5)  {Corr. RAG\\{\tiny Yan et al.}};
\node[below item] (agentic)  at (11.5,-1.5) {Agentic RAG\\{\tiny Singh et al.}};
\node[below item] (rarerag)  at (13.5,-1.5) {\textbf{RARE-RAG}\\{\tiny this work}};

% Connector lines (above)
\foreach \n/\x in {rag/0, dpr/0, hyde/4.6, selfrag/6.9, graphrag/9.2, raptor/9.2, chain/11.5, arag/13.5}{
  \draw[blue!40, line width=0.6pt] (\x, 0.25) -- (\n.south);
}
% Connector lines (below)
\foreach \n/\x in {colbert/0, beir/2.3, colbertv2/4.6, ragas/6.9, crag/9.2, agentic/11.5, rarerag/13.5}{
  \draw[teal!50, line width=0.6pt] (\x,-0.25) -- (\n.north);
}

\end{tikzpicture}
\caption{Timeline of key RAG-related publications (2020--2026).
  Items \emph{above} the axis are generation/architecture papers (blue);
  items \emph{below} are retrieval models and evaluation frameworks (teal).}
\label{fig:rag_timeline}
\end{figure}

% ------------------------------------------------------------
\section{From Language Models to Retrieval-Augmented Generation}
\label{sec:lm_to_rag}

\subsection{The Promise and Limits of Large Language Models}

Large language models trained on web-scale corpora have demonstrated remarkable
competence across a broad range of natural language tasks: they can translate, summarise,
answer factual questions, write and explain code, and reason through multi-step
problems with surprising reliability~\cite{ji2023hallucination}.
These capabilities arise from parametric knowledge --- patterns encoded into billions of
weights during training --- and they scale, roughly predictably, with both model size
and training data volume.

Yet the same training process that confers breadth imposes three fundamental limitations
on knowledge-intensive applications.
First, parametric knowledge has a fixed cutoff: anything published after the training
data collection ends is simply absent from the model.
For rapidly evolving fields such as pharmacology or clinical guidelines, a six-month
knowledge lag can be clinically significant.
Second, LLMs hallucinate.
The term covers a spectrum of failure modes --- confabulation of specific figures,
citation of non-existent papers, confident statement of incorrect facts --- but the
root cause is consistent: the model generates the most probable continuation of a
prompt, not the most faithful one~\cite{ji2023hallucination}.
On open-domain question-answering benchmarks, state-of-the-art models without retrieval
still produce incorrect answers in 20--40\% of cases depending on question difficulty.
Third, parametric models are opaque about their sources.
Even when correct, they cannot tell a user which document supports which claim, making
verification and trust calibration difficult.

\subsection{Non-Parametric Knowledge and the Retrieval Hypothesis}

A natural response to these limitations is to decouple knowledge storage from model
parameters.
Rather than encoding facts into weights at training time, store them in a searchable
index and retrieve the relevant subset at inference time.
This is the core of the \emph{retrieval hypothesis}: that language models can be
substantially more accurate and honest when conditioned on retrieved evidence,
even if that evidence is never seen during training.

The retrieval hypothesis has roots in open-domain question answering research from the
early 2000s, where systems such as IBM's Watson explicitly retrieved document passages
before generating answers.
What changed in 2020 was scale: the combination of dense vector retrieval with
billion-parameter language models produced a qualitatively different result from
earlier sparse-retrieval systems.

\subsection{Retrieval-Augmented Generation: The Original Framework}

Lewis \etal~\cite{lewis2020rag} formalised this combination as Retrieval-Augmented
Generation and demonstrated it on knowledge-intensive NLP tasks.
Their architecture is conceptually simple: given a query~$q$, retrieve the
top-$k$ documents $\{d_1, \ldots, d_k\}$ from a non-parametric memory using a dense
retriever, concatenate them with the query into a context window, and generate the
answer with a sequence-to-sequence model conditioned on that context.

The RAG framework introduced two key design choices that shaped subsequent work.
First, the use of a dense retriever (specifically DPR~\cite{karpukhin2020dpr}) rather
than BM25 enabled retrieval based on semantic similarity rather than lexical overlap,
which turned out to be critical for multi-hop and paraphrased queries.
Second, the end-to-end training setup allowed the retriever and generator to be
jointly optimised, though this joint training proved difficult to scale and was
largely abandoned in later work in favour of fixed retrievers paired with larger
generators.

On Natural Questions and TriviaQA, the RAG model outperformed parametric-only models
of comparable size by 5--10 points of Exact Match, establishing retrieval augmentation
as a reliable lever for factual accuracy.
The central insight --- that non-parametric memory can substitute for parametric
memorisation of facts --- has not been seriously challenged since.

% ------------------------------------------------------------
\section{Retrieval Foundations}
\label{sec:retrieval_foundations}

The quality of a RAG system's answers is bounded by the quality of its retrieval:
a generator conditioned on irrelevant or incorrect passages cannot recover the right
answer.
Understanding the retrieval layer is therefore essential to understanding why different
RAG architectures perform differently.

\subsection{Sparse Retrieval}

The oldest and most robust retrieval paradigm is sparse term matching.
BM25~\cite{robertson2009bm25}, a probabilistic extension of TF-IDF, scores document--query
relevance as a weighted sum of inverse document frequency terms,
with a saturation function that diminishes the impact of very high term frequencies.
Given a query $Q = \{q_1, \ldots, q_n\}$ and a document $D$, the BM25 score is:

\begin{equation}
  \text{BM25}(D, Q) = \sum_{i=1}^{n} \text{IDF}(q_i)
    \cdot \frac{f(q_i, D)\,(k_1 + 1)}{f(q_i, D) + k_1\!\left(1 - b + b\,\frac{|D|}{\text{avgdl}}\right)}
  \label{eq:bm25}
\end{equation}

\noindent where $f(q_i, D)$ is the frequency of term $q_i$ in document $D$,
$|D|$ is the document length, $\text{avgdl}$ is the average document length in the
corpus, and $k_1 \approx 1.5$, $b \approx 0.75$ are tunable parameters.
The $\text{IDF}$ term down-weights terms that appear in many documents.

Despite its simplicity, BM25 remains a strong baseline on many benchmarks: on the
BEIR heterogeneous retrieval benchmark~\cite{thakur2021beir}, BM25 outperforms
many neural retrievers on out-of-domain corpora.

SPLADE~v2 extended the sparse paradigm to learned sparse representations,
where a BERT encoder produces a sparse vector over the vocabulary rather than
a dense embedding.
This approach combines the efficiency of inverted-index retrieval with the
contextual understanding of transformers, and it achieves strong results on
in-domain benchmarks while maintaining interpretability through the explicit
vocabulary weights.

The principal limitation of sparse methods is lexical dependence: a query about
``myocardial infarction'' will fail to retrieve a highly relevant passage that
uses only ``heart attack''.
This brittleness motivated the dense retrieval line of work.

\subsection{Dense Retrieval}

Dense Passage Retrieval (DPR)~\cite{karpukhin2020dpr} encodes queries and passages
into fixed-dimensional vectors using separate BERT encoders, and retrieves by
approximate nearest-neighbour search in the embedding space.
Given a query encoder $E_Q$ and a passage encoder $E_P$, relevance is scored as

\begin{equation}
  \text{sim}(q, p) = E_Q(q)^\top E_P(p)
  \label{eq:dpr}
\end{equation}

\noindent and the top-$k$ passages are retrieved from a pre-built FAISS index using
maximum inner product search.
The encoders are fine-tuned on question--passage pairs so that relevant passages
cluster near their questions in the embedding space.

DPR demonstrated that dense retrieval can substantially outperform BM25 on
in-domain benchmarks such as Natural Questions, but the gains proved fragile:
DPR generalises poorly to new domains without fine-tuning.
Contriever~\cite{izacard2021contriever} addressed this by adopting a contrastive
self-supervised training objective that requires no labelled question--passage pairs,
producing an encoder that transfers better across domains.

The shared weakness of both DPR and Contriever is their single-vector representation:
each passage is compressed into one embedding, which inevitably loses detail.
For passages containing multiple facts, the embedding reflects their centroid ---
a problematic average that may not accurately represent any individual fact.

\subsection{Late Interaction: ColBERT}

ColBERT~\cite{khattab2020colbert} proposed a compromise between the efficiency of
single-vector retrieval and the expressiveness of cross-encoder reranking.
Rather than encoding a passage into a single vector, ColBERT produces a sequence
of token-level embeddings for both query ($\mathbf{q}_1,\ldots,\mathbf{q}_m$) and passage
($\mathbf{p}_1,\ldots,\mathbf{p}_n$).
Relevance is computed as the \emph{MaxSim} score:

\begin{equation}
  S(q, p) = \sum_{i=1}^{m}
    \max_{j=1}^{n}\; \mathbf{q}_i^\top \mathbf{p}_j
  \label{eq:maxsim}
\end{equation}

\noindent For each query token, find the maximum cosine similarity across all passage
token embeddings, then sum these maxima.
This late interaction mechanism is expressive enough to capture term-level alignment
while remaining amenable to pre-computation of passage embeddings.

ColBERTv2~\cite{santhanam2022colbertv2} improved the original design through
residual compression and distillation, reducing index size by an order of magnitude
while retaining most of the quality advantage.
ColBERT-style retrieval has since become a standard component of high-quality
RAG pipelines as both a first-stage retriever and a reranker.

\subsection{Hybrid Search and Cross-Encoder Reranking}

In practice, the strongest retrieval systems combine sparse and dense retrieval
through a fusion step.
Reciprocal Rank Fusion (RRF) is the most common approach: documents are ranked
separately by BM25 and a dense retriever, and the fused rank is computed as

\begin{equation}
  \text{RRF}(d) = \sum_{r \in \{\text{BM25},\, \text{dense}\}} \frac{1}{k + r(d)}
  \label{eq:rrf}
\end{equation}

\noindent where $r(d)$ is the rank of document $d$ in retriever $r$ and $k$ is a
smoothing constant (typically 60).
Alternatively, Weaviate and similar vector databases expose an $\alpha$ parameter
that linearly interpolates the BM25 and vector scores before ranking.
The \medrag{} system uses $\alpha = 0.5$ as the default, which was found empirically
to balance the strengths of both modalities on the HotpotQA development split
(Section~\ref{sec:hybrid_search}).

Cross-encoder rerankers take a different approach: rather than encoding query and
passage independently, they process the query--passage pair jointly with a
transformer, producing a relevance score that captures fine-grained interaction.
This is computationally expensive at scale but very accurate.
In practice, cross-encoders are applied only to the top-$k$ candidates from a
first-stage retriever, a pattern known as retrieve-then-rerank.
The \medrag{} system uses the BGE-reranker-v2-m3 model, a multilingual cross-encoder
from the BAAI general embedding family that performs competitively on the BEIR benchmark.

% ------------------------------------------------------------
\section{RAG Architecture Families}
\label{sec:rag_families}

The six years since Lewis \etal~\cite{lewis2020rag} have seen an explosion of
architectural variants.
Table~\ref{tab:rag_families} organises them into eight families, ordered by approximate
date of introduction.
The remainder of this section examines each family in turn.

\begin{table}[htbp]
  \centering
  \caption{Overview of surveyed RAG architecture families.
    Architectures are ordered chronologically.
    Benchmarks cited are those used in the original papers; direct comparison
    across rows is not valid without controlling for dataset and model size.}
  \label{tab:rag_families}
  \small
  \begin{tabularx}{\textwidth}{@{}llXX@{}}
    \toprule
    \textbf{Family} & \textbf{Year} & \textbf{Core idea} & \textbf{Key weakness} \\
    \midrule
    Vanilla RAG
      & 2020
      & Single-pass retrieve-then-generate
      & Fixed retrieval; no self-correction \\
    HyDE / Query Rewriting
      & 2022
      & Embed hypothetical answer for better retrieval alignment
      & Generator quality bottleneck; extra latency \\
    Self-RAG
      & 2023
      & Reflection tokens for adaptive retrieval and critique
      & Fine-tuning required; inference overhead \\
    Corrective RAG
      & 2024
      & Correctness evaluation with web-search fallback
      & Depends on web availability; hard to control \\
    Graph / Hierarchical RAG
      & 2024
      & Knowledge graphs or tree summaries over corpora
      & High indexing cost; graph construction is noisy \\
    Iterative Multi-hop RAG
      & 2025
      & Multi-step evidence chaining across retrieval calls
      & Latency scales with hop count; error propagation \\
    Agentic / Multi-Agent RAG
      & 2025
      & Planner--executor agents with tool-augmented retrieval
      & Complex orchestration; unpredictable cost \\
    \rarerag{} (this work)
      & 2026
      & Routing + hybrid retrieval + set-wise selection + grounding verifier
      & High latency on deep path; routing adds overhead \\
    \bottomrule
  \end{tabularx}
\end{table}

\subsection{Vanilla RAG}
\label{subsec:vanilla_rag}

The vanilla RAG pipeline, as introduced by Lewis \etal~\cite{lewis2020rag} and
subsequently simplified into a widely used pattern, operates in three stages.
First, the user's query is embedded by a retriever and used to retrieve the top-$k$
passages from an indexed corpus.
Second, the retrieved passages are concatenated with the query into a prompt.
Third, a language model generates the answer conditioned on that prompt.

The appeal of this pipeline is its simplicity: each component can be independently
swapped, and the only trainable interaction between them is through the prompt.
Figure~\ref{fig:vanilla_rag} illustrates the three-stage flow.

% ------------------------------------------------------------------
% Figure 2.2 — Vanilla RAG Pipeline
% ------------------------------------------------------------------
\begin{figure}[htbp]
\centering
\begin{tikzpicture}[
  font=\small,
  box/.style={
    rectangle, rounded corners=4pt,
    minimum width=2.5cm, minimum height=0.9cm,
    align=center, draw, thick
  },
  store/.style={
    cylinder, shape border rotate=90,
    aspect=0.25,
    minimum width=2.0cm, minimum height=0.9cm,
    align=center, draw, thick
  },
  arr/.style={-{Stealth[length=5pt]}, thick},
  lbl/.style={font=\footnotesize\itshape, midway, above=2pt}
]

% ---- nodes ----
\node[box, fill=blue!10,  draw=blue!50]   (query)   at (0,   0)    {Query $q$};
\node[box, fill=blue!15,  draw=blue!50]   (enc)     at (3.2, 0)    {Retriever\\$E_Q(q)$};
\node[store,fill=gray!10, draw=gray!50]   (index)   at (6.4, 0)    {Vector\\Index};
\node[box, fill=green!10, draw=green!50]  (topk)    at (6.4,-2.2)  {Top-$k$\\Passages};
\node[box, fill=orange!10,draw=orange!50] (prompt)  at (3.2,-2.2)  {Prompt\\$[q; d_1;\ldots;d_k]$};
\node[box, fill=red!8,    draw=red!40]    (llm)     at (0,  -2.2)  {LLM\\Generator};
\node[box, fill=blue!10,  draw=blue!50]   (answer)  at (0,  -4.0)  {Answer $a$};

% ---- arrows ----
\draw[arr] (query)  -- node[lbl]{embed} (enc);
\draw[arr] (enc)    -- node[lbl]{ANN search} (index);
\draw[arr] (index)  -- node[lbl,right=2pt]{retrieve} (topk);
\draw[arr] (topk)   -- node[lbl]{concatenate} (prompt);
\draw[arr] (prompt) -- node[lbl]{condition} (llm);
\draw[arr] (llm)    -- node[lbl,right=2pt]{generate} (answer);

\end{tikzpicture}
\caption{Vanilla RAG pipeline.
  The query is embedded by a retriever and used to retrieve the top-$k$ passages
  from a pre-built vector index.
  The passages are concatenated with the query into a prompt, which conditions
  the language model to generate the final answer.}
\label{fig:vanilla_rag}
\end{figure}

In practice, vanilla RAG with a capable retriever and a large generator achieves
competitive results on many benchmarks while remaining fast and predictable.
In our experiments (Chapter~\ref{cha:evaluation}), vanilla RAG achieves a Token~F1
of 0.537 and Exact Match of 0.316 on HotpotQA --- competitive with much more complex
architectures, at a fraction of their latency (9.1 seconds per query versus 20--31
seconds for enhanced modes).

The principal limitations of vanilla RAG are the lack of self-correction and the
fixed retrieval strategy.
If the top-$k$ documents are irrelevant or contradictory, the generator has no
mechanism to detect this and the answer quality degrades silently.
Subsequent architectures address precisely this failure mode.

\subsection{HyDE and Query Rewriting}
\label{subsec:hyde}

A recurring problem in RAG is the mismatch between query and document embeddings:
queries tend to be short and keyword-like, while passages are longer and more
elaborated, and the embeddings of the two may not be well aligned even after
retrieval-specific fine-tuning.

Gao \etal~\cite{gao2022hyde} proposed Hypothetical Document Embeddings (HyDE) as a
solution.
Rather than embedding the query directly, the system uses an LLM to generate a
\emph{hypothetical} document that would answer the query --- a plausible but
not necessarily accurate passage --- and embeds that hypothetical document instead.
The intuition is that the hypothetical document lives in the same embedding space as
real corpus passages, even if the facts it contains are hallucinated.
Retrieval then finds real passages that are semantically similar to the hypothetical
one.

HyDE is conceptually elegant and requires no additional training.
In our experiments, it achieves the second-highest faithfulness score (0.971) among all
evaluated modes, suggesting that the better-aligned retrieval does produce more
grounded context.
The cost is latency: HyDE requires an additional LLM call to generate the hypothetical
document before retrieval, pushing mean latency to 22.7 seconds per query.

Query rewriting pursues a related idea without the hypothetical document: a smaller
or cheaper model rewrites the user's query into one that is better suited for
retrieval, for instance by expanding acronyms, resolving coreferences, or
decomposing complex questions into simpler sub-questions.
Query rewriting is widely used in production RAG systems and can be combined with
any retrieval backbone.

\subsection{Self-RAG}
\label{subsec:self_rag}

Self-RAG~\cite{asai2023selfrag} represents a qualitative shift from the retrieve-then-
generate paradigm to a \emph{generate-with-retrieval-on-demand} paradigm.
The core contribution is a set of special reflection tokens injected into the
generation vocabulary:

\begin{itemize}
  \item \texttt{[Retrieve]} --- signals that additional retrieval is needed before
    continuing generation;
  \item \texttt{[IsRel]} / \texttt{[IsNotRel]} --- judges whether a retrieved passage
    is relevant to the query;
  \item \texttt{[IsSup]} / \texttt{[IsNotSup]} --- judges whether the generated
    statement is supported by the retrieved passage;
  \item \texttt{[IsUse]} --- rates the overall utility of a generated segment.
\end{itemize}

\noindent
A Self-RAG model is fine-tuned on data labelled with these tokens, learning when
to retrieve, how to assess retrieved evidence, and whether its own output is
grounded.
At inference time, the model interleaves retrieval calls with generation, inserting
passages where needed and backtracking when evidence is deemed insufficient.

Figure~\ref{fig:selfrag} illustrates the decision flow introduced by these tokens.

% ------------------------------------------------------------------
% Figure 2.3 — Self-RAG Decision Flow
% ------------------------------------------------------------------
\begin{figure}[htbp]
\centering
\begin{tikzpicture}[
  font=\small,
  startstop/.style={
    ellipse, draw, thick,
    minimum width=2.2cm, minimum height=0.8cm,
    fill=blue!10, draw=blue!50, align=center
  },
  process/.style={
    rectangle, rounded corners=3pt, draw, thick,
    minimum width=2.8cm, minimum height=0.8cm,
    fill=blue!8, draw=blue!40, align=center
  },
  decision/.style={
    diamond, draw, thick, aspect=2.2,
    minimum width=3.4cm, minimum height=0.8cm,
    fill=orange!10, draw=orange!50, align=center,
    inner sep=2pt
  },
  output/.style={
    rectangle, rounded corners=3pt, draw, thick,
    minimum width=2.8cm, minimum height=0.8cm,
    fill=green!10, draw=green!50, align=center
  },
  discard/.style={
    rectangle, rounded corners=3pt, draw, thick,
    minimum width=2.4cm, minimum height=0.7cm,
    fill=red!8, draw=red!30, align=center, font=\footnotesize
  },
  arr/.style={-{Stealth[length=5pt]}, thick},
  lbl/.style={font=\footnotesize, fill=white, inner sep=1pt}
]

% ---- nodes (column layout) ----
\node[startstop]  (q)       at (0,   0)    {Query $q$};
\node[decision]   (ret)     at (0,  -1.8)  {\texttt{[Retrieve?]}};
\node[process]    (retr)    at (3.5,-1.8)  {Retrieve\\passages};
\node[decision]   (isrel)   at (3.5,-3.5)  {\texttt{[IsRel?]}};
\node[discard]    (disc)    at (6.5,-3.5)  {Discard\\passage};
\node[process]    (gen)     at (0,  -5.2)  {Generate\\segment};
\node[decision]   (issup)   at (0,  -7.0)  {\texttt{[IsSup?]}};
\node[decision]   (isuse)   at (0,  -8.8)  {\texttt{[IsUse?]}};
\node[output]     (ans)     at (0, -10.5)  {Final answer $a$};
\node[discard]    (regen)   at (3.5,-7.0)  {Regenerate\\segment};

% ---- arrows ----
\draw[arr] (q)     -- (ret);
\draw[arr] (ret)   -- node[lbl, above]{yes} (retr);
\draw[arr] (ret)   -- node[lbl, right]{no} (gen);   % straight down
\draw[arr] (retr)  -- (isrel);
\draw[arr] (isrel) -- node[lbl, above]{not rel.} (disc);
\draw[arr] (isrel) -- node[lbl, right]{relevant} ++(0,-1.2)
              -| node[lbl]{} (gen);
\draw[arr] (gen)   -- (issup);
\draw[arr] (issup) -- node[lbl, above]{not sup.} (regen);
\draw[arr] (regen) -- ++(0, 0.8) -| (gen);   % loop back
\draw[arr] (issup) -- node[lbl, right]{supported} (isuse);
\draw[arr] (isuse) -- node[lbl, right]{useful} (ans);
\draw[arr] (isuse) -| node[lbl, above=4pt]{not useful} ++(3.2,0)
              |- (ret);   % loop back to top

\end{tikzpicture}
\caption{Self-RAG decision flow.
  Reflection tokens (\texttt{[Retrieve?]}, \texttt{[IsRel?]}, \texttt{[IsSup?]},
  \texttt{[IsUse?]}) are generated by the model itself at each step, controlling
  when to retrieve, which passages to keep, and whether to regenerate unsupported
  segments.}
\label{fig:selfrag}
\end{figure}

The results reported by Asai \etal show substantial gains over vanilla RAG on
fact verification and long-form generation tasks, particularly on questions that
require synthesising information from multiple sources.
In our evaluation, self-reflection achieves the highest Exact Match (0.329) of all
evaluated modes and the highest Token~F1 (0.545), confirming that explicit reasoning
about evidence quality does improve factual accuracy.
The cost is a 19.9-second mean latency and the requirement for a fine-tuned model ---
vanilla decoder models do not produce the reflection tokens needed for the mechanism
to operate.

\subsection{Corrective RAG}
\label{subsec:corrective_rag}

Corrective RAG (CRAG)~\cite{yan2024crag} takes the self-correction idea further by
evaluating the quality of retrieved documents before generation, rather than after.
A lightweight evaluator model scores each retrieved document on a spectrum from
\emph{correct} (relevant and accurate) to \emph{ambiguous} to \emph{incorrect}.
Based on these scores, the system chooses one of three actions:
if documents are correct, proceed to generation;
if ambiguous, refine the query and retrieve again;
if incorrect, fall back to a web search for fresh evidence.

The web-search fallback is the distinguishing feature: CRAG can recover from corpus
gaps that would silently mislead other architectures.
This makes it particularly suitable for questions about recent events or narrow
sub-domains not well represented in the static corpus.

In our HotpotQA experiments, corrective\_rag achieves a Token~F1 of 0.539 and
Exact Match of 0.323 at a mean latency of 11.8 seconds --- making it the best
cost-quality ratio among the enhanced architectures, roughly 2.5 times slower than
vanilla but significantly more accurate on the hardest questions.
However, the web-search fallback was effectively disabled in our benchmark
(no live web access), so these results represent the re-retrieval branch only.

\subsection{Adaptive RAG}
\label{subsec:adaptive_rag}

Adaptive-RAG~\cite{jeong2024adaptiverag} proposes classifying queries by complexity
before retrieval, routing them to different pipeline variants: no retrieval for simple
factoid questions answerable from parametric memory, single-step retrieval for
moderate queries, and multi-step retrieval for complex multi-hop questions.
The classifier is a small model fine-tuned on query complexity labels derived from
benchmark datasets.

The key insight is that not all queries benefit equally from retrieval, and applying
a heavyweight multi-step pipeline to simple questions wastes compute without improving
quality.
RARE-RAG (Chapter~\ref{cha:rare_rag}) builds on this routing idea but extends it with
domain-specific risk classification and integrates it with set-wise evidence selection
and a grounding verifier.

\subsection{Graph and Hierarchical RAG}
\label{subsec:graph_rag}

Standard RAG retrieves flat passages.
For corpora where important information is distributed across many documents and
latent connections between entities matter, flat retrieval misses structure that
a human expert would use.
Graph and hierarchical approaches address this by building richer representations
of the corpus at indexing time.

GraphRAG~\cite{edge2024graphrag}, developed by Microsoft Research, constructs a
knowledge graph from the corpus using an LLM-driven entity and relationship
extraction pipeline.
At query time, it performs graph traversal to find relevant communities of entities,
then generates summaries of those communities as retrieval context.
GraphRAG demonstrated substantial improvements over flat RAG on global summarisation
tasks --- questions that require synthesising information spread across an entire
corpus rather than retrieving specific passages.

RAPTOR~\cite{sarthi2024raptor} takes a complementary approach: rather than building
a graph, it recursively summarises clusters of related documents into a tree
structure, then retrieves from both leaf nodes (raw passages) and internal nodes
(summaries).
This allows the retriever to operate at multiple levels of abstraction, answering
broad questions from summaries and narrow questions from raw passages.

Both approaches share a significant drawback: the indexing pipeline is expensive
and must be re-run when the corpus changes.
For dynamic medical corpora (where new drug interaction data is published regularly),
this imposes a meaningful operational burden.

\subsection{Iterative Multi-Hop RAG}
\label{subsec:multihop}

Multi-hop questions require chaining evidence across multiple documents: ``What drug
is contraindicated with the anticoagulant prescribed for the condition that causes
blood clots?'' cannot be answered from any single passage.
The iterative multi-hop family addresses this by running multiple retrieval calls in
sequence, where each call uses the evidence collected so far to formulate the next
sub-query.

Chain-of-Retrieval Augmented Generation~\cite{wang2025chain} implements this as a
structured generation process where the model produces interleaved reasoning steps
and retrieval queries, executing each retrieval call before continuing.
This is similar in spirit to Self-RAG but focuses specifically on the multi-hop
evidence chain rather than general retrieval quality assessment.

In our experiments, iterative\_multihop achieves a Token~F1 of 0.533 and Exact
Match of 0.314 at a mean latency of 26.0 seconds.
The latency cost is substantial because each hop requires a full retrieval call,
and HotpotQA questions often require two to three hops.
The context recall (0.953) is also the second-lowest among all modes, suggesting that
the iterative process does not always recover evidence that a single broad retrieval
would have found.

\subsection{Agentic and Multi-Agent RAG}
\label{subsec:agentic_rag}

The most recent and most flexible architectural family treats RAG as an agent
planning problem.
Rather than a fixed pipeline, an agentic RAG system consists of one or more LLM
agents that decide at each step whether to retrieve, which query to issue, whether
the collected evidence is sufficient, and whether to involve specialised sub-agents
for particular sub-tasks.

Singh \etal~\cite{singh2025agenticrag} survey this emerging landscape, organising
it by number of agents (single vs.\ multi-agent), control structure (sequential,
hierarchical, collaborative), and autonomy level.
At the high end, systems such as A-RAG~\cite{du2026arag} expose three distinct
retrieval interfaces --- keyword search, semantic search, and chunk read --- directly
to the LLM as callable tools, allowing the model to compose retrieval strategies
rather than following a predetermined pipeline.

MA-RAG~\cite{nguyen2025marag} decomposes the pipeline into specialised roles:
a planner that decomposes the question, a step definer that specifies each retrieval
action, an extractor that processes retrieved passages, and a QA agent that synthesises
the final answer.
This division of labour reduces the cognitive load on any single model and allows
different components to be optimised independently.

MADAM-RAG~\cite{wang2025madamrag} explores a different angle: multi-agent debate
for questions with ambiguous or conflicting evidence.
Multiple agents retrieve evidence independently and argue for different candidate
answers, with a judge agent synthesising the debate into a final response.
While compelling for disambiguation tasks, this design proved poorly suited to
factoid multi-hop QA in our experiments: MADAM-RAG achieves the lowest Token~F1
(0.487) and Exact Match (0.282) of all evaluated modes, at the highest latency (29.0
seconds) after \rarerag{}.
The debate mechanism adds overhead without benefit when the question has a unique,
retrievable answer.

In our evaluation, the multi\_agent mode (a single-orchestrator multi-tool agent)
achieves a surprising result: Token~F1 of 0.525 and mean latency of only 10.2 seconds
--- the second-fastest mode after vanilla, due to parallel tool calls reducing
wall-clock time despite higher total token count.
This suggests that well-designed agentic architectures can be both capable and
efficient when parallelism is exploited.

% ------------------------------------------------------------
\section{Evaluation of RAG Systems}
\label{sec:rag_evaluation}

Evaluating RAG systems is harder than evaluating either retrieval or generation
in isolation, because quality depends on the interaction between the two components.
A system that retrieves perfect evidence but generates a garbled answer, or one that
generates a fluent response from irrelevant passages, will both fail --- in different
ways that aggregate metrics may obscure.
This section reviews the metrics and benchmarks used in this thesis.

\subsection{Lexical Overlap Metrics}

The earliest and most widely used answer-quality metrics are lexical overlap measures.

\textbf{Exact Match (EM)} scores a prediction as correct only if it matches the
gold answer string exactly after normalisation (lowercasing, removing articles and
punctuation).
EM is strict: a correct but differently phrased answer scores zero.

\textbf{Token F1} relaxes this by computing a token-level F1 score between
prediction and gold answer, treating both as bags of tokens.
Token~F1 is more forgiving of paraphrase but still penalises missing content words.

\textbf{ROUGE-L}~\cite{lin2004rouge} measures the longest common subsequence between
prediction and reference, normalised by the lengths of both.
ROUGE-L is commonly used for longer answer evaluation.

All three metrics have the well-known limitation that they measure surface form rather
than semantic correctness: a prediction that says ``approximately 47 milligrams'' when
the gold answer is ``47~mg'' will score poorly despite being correct.
For medical advisory, where the relevant answers are often drug names, doses, or
interaction categories with specific canonical forms, this limitation is less
acute than in open-ended generation tasks.

\subsection{The RAGAS Framework}

The RAGAS framework~\cite{es2023ragas} addresses the limitations of lexical metrics
by introducing four metrics that capture different dimensions of RAG quality:

\textbf{Faithfulness} measures the fraction of claims in the generated answer that
are supported by the retrieved context.
Claims $C = \{c_1, \ldots, c_m\}$ are extracted from the answer by an LLM,
and each claim is checked against the context by the same LLM acting as a natural
language inference judge.
Formally:
\begin{equation}
  \text{Faithfulness}(a, \mathcal{C}) =
    \frac{\bigl|\{c_i \in C : c_i \text{ is supported by } \mathcal{C}\}\bigr|}{|C|}
  \label{eq:faithfulness}
\end{equation}
where $\mathcal{C} = \{d_1,\ldots,d_k\}$ is the set of retrieved passages.
A score of 1.0 means every claim in the answer can be traced to a specific passage;
a score below 0.8 suggests significant hallucination.
Faithfulness is the metric most directly relevant to trustable medical advisory:
an answer that contradicts its own supporting documents is dangerous regardless of
whether it happens to be correct.

\textbf{Answer Relevance} measures how closely the generated answer addresses the
original question.
It is computed by generating $n$ reverse questions $\hat{q}_1,\ldots,\hat{q}_n$
from the answer using an LLM and measuring their average cosine similarity to the
original query:
\begin{equation}
  \text{AnswerRelevance}(a, q) =
    \frac{1}{n} \sum_{i=1}^{n}
    \frac{E(\hat{q}_i) \cdot E(q)}{\|E(\hat{q}_i)\|\,\|E(q)\|}
  \label{eq:answerrel}
\end{equation}
where $E(\cdot)$ denotes the sentence embedding.
A high-relevance answer stays on-topic; a low-relevance one drifts into tangential
information.

\textbf{Context Recall} measures how much of the information in the gold answer
is covered by the retrieved context.
It is the only RAGAS metric that requires a ground-truth answer, and it captures
retrieval completeness rather than answer quality.

\textbf{Context Precision} measures the fraction of retrieved chunks that are
actually useful for answering the question.
It captures retrieval noise: a system that retrieves many irrelevant passages alongside
the relevant ones will score low on context precision.

RAGAS uses an LLM as a judge for faithfulness and answer relevance, which introduces
a dependency: systems evaluated with the same model family used for generation may
benefit from a systematic bias.
This limitation is acknowledged in Section~\ref{sec:limitations}.

\subsection{Benchmarks and Datasets}
\label{subsec:benchmarks}

\textbf{HotpotQA}~\cite{yang2018hotpotqa} is a multi-hop question-answering dataset
based on English Wikipedia.
Each question requires synthesising information from exactly two supporting documents
to arrive at the answer.
The dataset contains approximately 113,000 question--answer pairs; we use a
stratified sample of 1,000 questions per mode in our evaluation.
HotpotQA is well-suited for RAG evaluation because the multi-hop structure stresses
the retrieval component: a system that retrieves only one of the two required documents
cannot answer correctly.

\textbf{BEIR}~\cite{thakur2021beir} is a heterogeneous retrieval benchmark covering
18 datasets from diverse domains, designed to evaluate retrieval generalisation.
It is not used directly in this thesis but is referenced when discussing the
relative strengths of retrieval models.

% ------------------------------------------------------------
\section{Related Work in Clinical and Medical NLP}
\label{sec:medical_nlp}

The application of natural language processing to clinical medicine predates the LLM
era by several decades, with rule-based systems for medical coding and
information extraction dating to the 1990s.
This section focuses on the more recent literature directly relevant to drug
interaction advisory.

\subsection{Clinical Decision Support Systems}

Clinical decision support systems (CDSS) assist clinicians by surfacing relevant
guidelines, flagging potential drug interactions, and suggesting diagnostic
possibilities based on patient data.
Rule-based CDSS such as the First DataBank and Micromedex databases encode drug
interaction knowledge as manually curated rule sets with severity classifications.
These systems are effective within their curated scope but cannot generalise to
combinations or dosing regimens not explicitly enumerated.

Machine learning approaches have been explored for automated DDI prediction from
molecular structure and biological pathway data, with graph neural networks on
drug--target interaction networks achieving strong results on benchmark datasets.
These approaches, however, operate on structured chemical data rather than natural
language clinical queries, and they do not provide explanations accessible to
clinicians without specialised training.

\subsection{LLMs for Drug Interaction Detection}

Several groups have evaluated LLMs on drug interaction questions, with mixed results.
Properly prompted GPT-4 achieves reasonable accuracy on standardised DDI
question-answering tasks from DrugBank, but the accuracy is inconsistent and
the model frequently presents its answers with false confidence.
The hallucination problem is particularly acute for DDI: an LLM that
confabulates an interaction between two drugs that are actually safe in combination
could cause unnecessary treatment changes, while missing a real interaction could
cause harm.

RAG has been proposed as a mitigation: by grounding DDI answers in retrieved passages
from authoritative sources (FDA drug labels, clinical pharmacology textbooks), the
system's claims become traceable and its failures become detectable.
The \medrag{} system implements this approach as a technology demonstrator,
using the OpenFDA drug label corpus as the primary knowledge base for drug
interaction queries.

\subsection{Trustability and Explainability}

Trustability in clinical AI has two components: the system must be accurate enough
to be useful, and it must be transparent enough to allow clinicians to calibrate
their reliance on it.
The second requirement is often neglected in ML research but is central to real-world
clinical deployment~\cite{ji2023hallucination}.

RAG contributes to trustability through citation: every answer can be linked to the
specific passages that support it, allowing a clinician to verify the claim directly.
The \rarerag{} architecture (Chapter~\ref{cha:rare_rag}) takes this further with an
explicit grounding verifier that checks whether the generated answer is supported by
the retrieved context, and abstains from answering when it cannot verify the claims.
This abstention mechanism is particularly important in the DDI domain, where a
``I cannot reliably answer this question based on available evidence'' is safer
than a plausible-sounding but unverified response.

% ------------------------------------------------------------------
% Figure 2.4 — Architecture Comparison: Retrieval Adaptivity vs. Generation Complexity
% ------------------------------------------------------------------
\begin{figure}[htbp]
\centering
\begin{tikzpicture}[
  font=\small,
  dot/.style={circle, draw, thick, minimum size=0.55cm, inner sep=1pt, align=center},
  arr/.style={-{Stealth[length=4pt]}, thick, gray!60}
]

\draw[arr] (-0.3, 0) -- (10.5, 0)
  node[right, font=\small\bfseries, black] {Retrieval adaptivity};
\draw[arr] (0, -0.3) -- (0, 8.2)
  node[above, font=\small\bfseries, black] {Generation complexity};

\node[font=\footnotesize, gray] at (1.5, -0.45) {fixed};
\node[font=\footnotesize, gray] at (5.0, -0.45) {moderate};
\node[font=\footnotesize, gray] at (9.0, -0.45) {fully adaptive};
\node[font=\footnotesize, gray, rotate=90] at (-0.65, 1.5) {single-pass};
\node[font=\footnotesize, gray, rotate=90] at (-0.65, 4.5) {iterative};
\node[font=\footnotesize, gray, rotate=90] at (-0.65, 7.2) {multi-agent};

\foreach \x in {2.5, 5.0, 7.5}{
  \draw[gray!15, dashed] (\x, 0) -- (\x, 8.0);
}
\foreach \y in {2.5, 5.0, 7.5}{
  \draw[gray!15, dashed] (0, \y) -- (10.2, \y);
}

\node[dot, fill=blue!20,   draw=blue!60]   (van)   at (1.2, 1.2) {};
\node[font=\footnotesize, above=3pt]  at (van)   {Vanilla RAG};

\node[dot, fill=cyan!20,   draw=cyan!60]   (hyde)  at (3.0, 1.5) {};
\node[font=\footnotesize, above=3pt]  at (hyde)  {HyDE};

\node[dot, fill=cyan!25,   draw=cyan!60]   (crag)  at (5.5, 2.0) {};
\node[font=\footnotesize, above=3pt]  at (crag)  {Corr. RAG};

\node[dot, fill=orange!25, draw=orange!60] (self)  at (5.0, 3.8) {};
\node[font=\footnotesize, right=3pt]  at (self)  {Self-RAG};

\node[dot, fill=purple!20, draw=purple!50] (graph) at (3.5, 4.5) {};
\node[font=\footnotesize, left=3pt]   at (graph) {Graph RAG};

\node[dot, fill=teal!20,   draw=teal!60]   (iter)  at (6.0, 5.0) {};
\node[font=\footnotesize, right=3pt]  at (iter)  {Iter. Multi-hop};

\node[dot, fill=red!15,    draw=red!50]    (madam) at (7.5, 6.5) {};
\node[font=\footnotesize, above=3pt]  at (madam) {MADAM-RAG};

\node[dot, fill=green!25,  draw=green!60,
      minimum size=0.78cm, line width=1.5pt] (rare) at (8.8, 5.8) {};
\node[font=\footnotesize\bfseries, above=3pt] at (rare) {\rarerag{}};

\draw[green!50!black, thick, dashed, rounded corners=6pt]
  (1.2,1.2) .. controls (3.8,2.2) and (6.2,4.0) .. (8.8,5.8);
\node[font=\footnotesize, green!60!black, rotate=26] at (5.0, 3.2)
  {Pareto frontier};

\end{tikzpicture}
\caption{Positioning of the eight surveyed RAG families along two design dimensions.
  \emph{Retrieval adaptivity} captures how dynamically the system selects its
  retrieval strategy; \emph{generation complexity} captures the number of
  reasoning or generation steps.
  The dashed curve approximates the quality--complexity Pareto frontier;
  \rarerag{} is positioned at its high-adaptivity end.}
\label{fig:arch_comparison}
\end{figure}

% ------------------------------------------------------------
\section{Summary}
\label{sec:background_summary}

This chapter has traced the development of RAG from its foundations in dense
retrieval and the original Lewis \etal framework, through six families of increasingly
sophisticated architectures, to the evaluation methodology used in this thesis.

Several observations motivate the \rarerag{} design presented in
Chapter~\ref{cha:rare_rag}.
First, no existing architecture in the literature routes queries by risk or clinical
significance before selecting a retrieval strategy: Adaptive-RAG routes by linguistic
complexity but not domain-specific risk.
Second, all surveyed architectures rank retrieved evidence individually: set-wise
evaluation of evidence complementarity is absent from the standard RAG toolbox.
Third, abstention mechanisms --- declining to answer when evidence is insufficient ---
are treated as optional features rather than first-class design goals, despite their
importance in high-stakes advisory settings.
\rarerag{} addresses all three gaps in a unified architecture.
